\documentclass[12pt]{article}
\usepackage[T1]{fontenc}
\usepackage[utf8]{inputenc}
\usepackage{lmodern}
\usepackage{textcomp}
\usepackage{enumitem}
\usepackage{geometry}
\geometry{margin=1in}
\begin{document}
\begin{center}
Answer Key - Sociology - Undergraduate Level Assessment 11 \
\small (Answers are ordered exactly as in the question paper.)
\end{center}

\section*{Multiple Choice}
\begin{enumerate}[label=\arabic*.]
\item \textbf{Answer:} D
\\ \textit{Options:} A) they had rediscovered themselves as 'new men'; B) their wives were at home and nagged them all the time; C) exciting gadgets like the hoover and electric iron were invented; D) they were unemployed or both partners worked full time
\\ \textit{Note:} Sullivan's study found that men engaged in a greater share of housework when they were unemployed or when both partners worked full time, as these situations often necessitate a more equitable division of domestic responsibilities to balance out the demands of their employment status.
\item \textbf{Answer:} D
\\ \textit{Options:} A) including union representatives in processes of policy decision making; B) creating links between the state and corporate organizations; C) recruiting women into full time paid employment; D) including working class organizations in political bargaining and representation
\\ \textit{Note:} The correct answer is justified because 'incorporation' in the context of the nineteenth-century labour movement refers to the process of formally integrating working-class organizations into the political system, allowing them to participate in negotiations and influence policymaking on labor issues.
\item \textbf{Answer:} A
\\ \textit{Options:} A) ways of acting, thinking, and feeling that are collective and social in origin; B) the way scientists construct knowledge in a social context; C) data collected about social phenomena that are proven to be correct; D) ideas and theories that have no basis in the external, physical world
\\ \textit{Note:} The correct answer reflects Durkheim's definition of social facts as phenomena that exist outside the individual and shape collective behavior, emphasizing their nature as products of social contexts rather than individual actions.
\item \textbf{Answer:} C
\\ \textit{Options:} A) a class in itself; B) a class by itself; C) a class for itself; D) a ruling class
\\ \textit{Note:} Marx argued that as workers united through factory-based production, they would gain class consciousness and recognize their shared interests and struggles, thus evolving into a "class for itself" that actively advocates for their rights and collective identity.
\item \textbf{Answer:} A
\\ \textit{Options:} A) human behaviour is meaningful, and varies between individuals and cultures; B) it is difficult for sociologists to gain access to a research laboratory; C) sociologists are not rational or critical enough in their approach; D) we cannot collect empirical data about social life
\\ \textit{Note:} The correct answer highlights that human behavior is influenced by individual experiences, cultural contexts, and social meanings, making it inherently more complex and variable than phenomena in the natural world, which follow consistent, universal laws.
\end{enumerate}

\section*{True / False}
\begin{enumerate}[label=\arabic*.]
\setcounter{enumi}{5}
\item \textbf{Answer:} False
\\ \textit{Options:} A) True; B) False
\\ \textit{Note:} The experimental method relies on systematic observation and empirical evidence rather than personal beliefs and values to guide research decisions.
\item \textbf{Answer:} True
\\ \textit{Options:} A) True; B) False
\\ \textit{Note:} Totalitarian societies restrict freedom of movement to prevent dissent, monitor the population, and ensure that citizens remain under the state's control, thereby maintaining their authority and power.
\item \textbf{Answer:} True
\\ \textit{Options:} A) True; B) False
\\ \textit{Note:} Fulcher \& Scott argue that a community can be defined by social relationships and shared interests rather than merely by a specific geographical location, emphasizing the importance of social connections and interactions over physical space.
\item \textbf{Answer:} True
\\ \textit{Options:} A) True; B) False
\\ \textit{Note:} The human relations approach is based on the belief that fostering positive interpersonal relationships, enhancing communication, and prioritizing employee well-being lead to improved organizational performance and job satisfaction.
\item \textbf{Answer:} False
\\ \textit{Options:} A) True; B) False
\\ \textit{Note:} Sociology employs systematic data collection and empirical evidence through various research methods, such as surveys and experiments, to analyze social phenomena objectively rather than solely relying on subjective interpretations.
\end{enumerate}

\section*{Long Form Response}
\begin{enumerate}[label=\arabic*.]
\setcounter{enumi}{10}
\item \textbf{Answer:} The "mollies" were a deviant subculture of homosexuals in seventeenth and eighteenth century London, characterized by their distinct social practices, fashion choices, and secretive gatherings. Often associated with a flamboyant lifestyle, mollies would frequent specific venues, such as taverns and private homes, where they could express their identities away from societal norms. This subculture highlighted the tension between emerging modern identities and traditional societal expectations, as the mollies' existence challenged contemporary views on sexuality and gender roles. The societal implications were significant, as their presence prompted discussions about morality, legality, and the nature of deviance, ultimately contributing to the early discourse on sexual orientation and identity in Western society. Thus, the mollies played a crucial role in shaping the understanding of homosexuality in a period marked by rigid social hierarchies and norms.
\\ \textit{Note:} correct because it accurately identifies the characteristics of the mollies, such as their flamboyant lifestyle and social practices in secretive venues, while also addressing the broader societal implications of their existence in the context of conflicts between emerging modern identities and traditional norms in seventeenth and eighteenth century London.
\item \textbf{Answer:} Globalization is characterized by the increasing interconnectedness of economies, cultures, and populations across the globe, driven by advancements in technology, trade liberalization, and the movement of information. This phenomenon challenges the traditional power of nation-states by undermining their ability to control economic policies and cultural narratives within their borders, as multinational corporations and international organizations often exert greater influence. As transnational issues such as climate change and migration emerge, the responsibilities of nation-states are increasingly shared or complicated by global governance structures, thus diluting their sovereignty. However, while globalization may weaken the authority of individual states, it also compels them to adapt and extend their power by engaging in new forms of diplomacy and collaboration, highlighting a complex interplay between global forces and national interests.
\\ \textit{Note:} The answer accurately identifies globalization's key features, such as economic interconnectedness and cultural exchange, while highlighting how these elements diminish the sovereignty and control of nation-states over their own policies and narratives.
\end{enumerate}

\vfill
\noindent\textit{End of Answer Key}
\end{document}